% Part: normal-modal-logic
% Chapter: tableaux
% Section: introduction

\documentclass[../../../include/open-logic-section]{subfiles}

\begin{document}

\olfileid{nml}{tab}{int}

\olsection{引言}

分析表是带号公式的某种（向下分支的）树，所谓带号公式即由真值符号（$\True$ 或
$\False$）与一个语句
\[
\sFmla{\True}{!A} \text{ 或 } \sFmla{\False}{!A}.
\]
组成的对。分析表由若干\emph{假设}开始。其余的每个
带号公式都是通过应用某条推理规则生成的。有些推理规则向树的某个
树梢添加一个或多个带号公式；另一些规则则添加两个新树梢，从而产生两个分支。
规则生成的带号公式，其公式比应用该规则的带号公式的公式
复杂度更低。当一个分支同时包含 $\sFmla{\True}{!A}$ 和
$\sFmla{\False}{!A}$ 时，我们称该分支是\emph{封闭的}。如果分析表中的每个
分支都封闭，那么整个分析表就是封闭的。一个封闭的分析表构成了一个推导，
它表明作为分析表起始假设的那组带号公式是
不可满足的。这可以用来定义 $\Proves$ 关系：
$\Gamma \Proves !A$ 当且仅当存在某个有限子集~$\Gamma_0 = \{!B_1,
\dots, !B_n\} \subseteq \Gamma$，使得以假设
\[
\{\sFmla{\False}{!A}, \sFmla{\True}{!B_1}, \dots, \sFmla{\True}{!B_n}\}.
\]
为起点存在一个封闭的分析表。

对于模态逻辑，我们既要扩展带号公式的概念，又要添加覆盖~\iftag{prvBox}{$\Box$\iftag{prvDiamond}{ 和
    $\Diamond$}}{$\Diamond$} 的规则。在模态分析表中，公式除了具有真值符号（$\True$ 或
$\False$）之外，还具有\emph{前缀}~$\sigma$。前缀是正整数的非空序列，即 $\sigma \in (\PosInt)^* \setminus
\{\emptyseq\}$。书写这些前缀时，我们省略外围的
$\tuple{\ }$，并用~$.$ 代替 $,$ 来分隔各个元素。如果 $\sigma$ 是一个前缀，那么 $\sigma.n$ 就是 $\sigma \concat \tuple{n}$；例如，若 $\sigma = 1.2.1$，则 $\sigma.3$ 就是
$1.2.1.3$。因此，
\[
\sFmla{\True}{\Box !A \lif !A}[1.2]
\]
就是一个\emph{带前缀的带号公式}（或简称为\emph{带前缀公式}）。

直观上，前缀命名了模型中的一个世界，该世界可能满足分析表某分支上的公式；如果 $\sigma$ 命名了
某个世界，那么 $\sigma.n$ 就命名了一个从（由~$\sigma$ 命名的）世界可及的世界。

\end{document}